%%%%%%%%%%%%%%%%%%%%%%%%%%%%%%%%%%%%%
% Samy Tamim — CV (LaTeX), version française senior/ATS
% Compilateur : pdflatex / XeLaTeX
% Version FR de resume/resume.tex — garder les deux en synchronisation.
%%%%%%%%%%%%%%%%%%%%%%%%%%%%%%%%%%%%%
\documentclass[10pt,a4paper]{article}

% ── Packages ─────────────────────────────────────────────────────────────────
\usepackage[top=0.50in, bottom=0.50in, left=0.55in, right=0.55in]{geometry}
\usepackage{xcolor}
\usepackage{hyperref}
\usepackage{enumitem}
\usepackage{titlesec}
\usepackage{array}
\usepackage{tabularx}

% ── Fonts (Carlito, métriquement compatible avec Calibri) ────────────────────
\newif\ifluafonts
\ifdefined\XeTeXrevision\luafontstrue\fi
\ifdefined\directlua\luafontstrue\fi

\ifluafonts
  \usepackage{fontspec}
  \usepackage{textcomp}
  \usepackage[expansion=false]{microtype}
  \IfFontExistsTF{Carlito}
    {\setmainfont{Carlito}[Ligatures=TeX, Numbers=Lining, BoldFont={Carlito Bold}]}
    {\setmainfont{Latin Modern Roman}[Ligatures=TeX]%
     \PackageWarning{resume}{Carlito system font not found - falling back. Install fonts-crosextra-carlito for Carlito}}
  \newcommand{\skill}[1]{\textbf{#1}}
\else
  \usepackage[T1]{fontenc}
  \usepackage[utf8]{inputenc}
  \usepackage{textcomp}
  \IfFileExists{carlito.sty}
    {\usepackage[lining]{carlito}%
     \renewcommand{\familydefault}{\sfdefault}}
    {\usepackage{lmodern}%
     \PackageWarning{resume}{carlito.sty not found - falling back to lmodern. Install texlive-fonts-extra for Carlito}}
  \usepackage[expansion=false]{microtype}
  \DisableLigatures{family=*}
  \newcommand{\skill}[1]{\textbf{#1}}
\fi

% ── Couleurs ─────────────────────────────────────────────────────────────────
\definecolor{accent}{HTML}{1A1A2E}
\definecolor{link}{HTML}{2563EB}

% ── Hyperref ─────────────────────────────────────────────────────────────────
\hypersetup{colorlinks=true, urlcolor=link, hidelinks, pdfnewwindow=true, linkbordercolor={1 1 1},
  pdftitle={Samy Tamim - CV Développeur Backend et IA},
  pdfauthor={Samy Tamim},
  pdfsubject={Développeur Backend et IA - Systèmes Distribués, Architecture IoT et Agents IA},
  pdfkeywords={CV, Développeur Backend, Développeur IA, Systèmes Distribués, Microservices, Architecture Événementielle, Kubernetes, Kafka, MQTT, Go, Python, Java, TypeScript, Jumeau Numérique, Maintenance Prédictive, Agents IA, MCP, LLM, CI/CD, RNCP Niveau 7, Bac+5}}

% ── Style des sections ───────────────────────────────────────────────────────
\titleformat{\section}
  {\normalsize\bfseries\color{accent}}
  {}{0em}{}
  [\vspace{-0.38em}\rule{\linewidth}{0.7pt}\vspace{-0.33em}]

\titlespacing{\section}{0pt}{0.45em}{0.22em}

% ── Style des listes ─────────────────────────────────────────────────────────
\setlist[itemize]{
  leftmargin = 1.1em,
  itemsep    = 0.90pt,
  topsep     = 1.20pt,
  parsep     = 0pt,
  partopsep  = 0pt
}

% ── Commandes utilitaires ────────────────────────────────────────────────────
\newcommand{\entry}[2]{%
  \noindent\textbf{#1}\hfill{\small\textbf{#2}}\vspace{0.04em}%
}
\newcommand{\subentry}[2]{%
  \noindent{\small #1}\hfill{\footnotesize #2}\vspace{0.04em}%
}

\pagestyle{empty}
\setlength{\parindent}{0pt}
\setlength{\parskip}{0pt}
% French words can't hyphenate (no babel/french), so give TeX extra stretch to
% avoid overfull justified lines in long words like "événementielle".
\setlength{\emergencystretch}{1.5em}
\renewcommand{\baselinestretch}{0.96}

% ─────────────────────────────────────────────────────────────────────────────
\begin{document}

% ── En-tête ──────────────────────────────────────────────────────────────────
\begin{center}
  {\LARGE\bfseries\color{accent} Samy Tamim}\\[1.5pt]
  {\small Développeur Backend \& IA}\\[1.5pt]
  \footnotesize
  +33 781 805 041 \enspace\textbullet\enspace \href{mailto:contact@samytamim.com}{contact@samytamim.com} \enspace\textbullet\enspace \href{https://linkedin.com/in/samy-tamim}{https://linkedin.com/in/samy-tamim} \enspace\textbullet\enspace \href{https://samytamim.com}{https://samytamim.com} \enspace\textbullet\enspace \href{https://github.com/mylastresort}{https://github.com/mylastresort}
\end{center}

% ── Résumé ───────────────────────────────────────────────────────────────────
\section{Résumé professionnel}
Développeur backend et IA avec plus de 2 ans d'expérience, spé\-cialisé dans les systèmes distribués, les microservices et l'architecture événementielle pour des plateformes IoT et temps réel, y compris des agents LLM. Dirige une plateforme de jumeau numérique avec maintenance prédictive~; pratique approfondie de \skill{Go}, \skill{Java}, \skill{Python}, \skill{TypeScript}, \skill{Kubernetes}, \skill{Kafka} et \skill{MQTT}. Axé sur scalabilité, tolérance aux pannes, sécurité et pratiques d'ingénierie (revue de code, TDD, CI/CD).

% ── Expérience ───────────────────────────────────────────────────────────────
\section{Expérience professionnelle}

\entry{Ingénieur Logiciel}{Juin 2024 – Déc. 2026}\\
\subentry{\href{https://um6p.ma/}{UM6P - Université Mohammed VI Polytechnique}}{Rabat, Maroc}
\begin{itemize}
  \item Conception et direction de l'architecture d'une plateforme de jumeau numérique pour la supervision IoT et la maintenance prédictive, ingérant les données de capteurs d'équipements industriels via \skill{MQTT} avec analyses en temps réel.
  \item Développement d'un assistant d'exploitation IoT multi-agents sur \skill{Google ADK}~: sous-agents spécialisés (finder d'actifs, agent de télémétrie, agent PdM) en agent-as-tool, modèles \skill{Ollama} locaux et \skill{LiteLLM} en couche proxy pour des appels d'outils fiables, streaming et raisonnement, et environ 98 outils exposés via \skill{MCP} avec garde-fous contre les hallucinations~; évaluation des modèles et des agents via des benchmarks fondés sur les trajectoires.
  \item Développement de microservices \skill{Kafka} et \skill{Python} avec cache \skill{Redis}, architecture entièrement \skill{sans état} (stateless), y compris les consommateurs et les flux.
  \item Évaluation de plusieurs \skill{SPOF} de la plateforme pour garantir la résilience du système.
\end{itemize}

\vspace{0.40em}

\entry{Stagiaire Développement Full-Stack}{Déc. 2023 – Mai 2024}\\
\subentry{\href{https://um6p.ma/}{UM6P - Université Mohammed VI Polytechnique}}{Rabat, Maroc}
\begin{itemize}
  \item Développement, en équipe \skill{Agile} pluridisciplinaire, d'une plateforme suivant plus de 200 mouvements de navires par jour pour le commerce d'engrais, avec informations temps réel sur la chaîne d'approvisionnement et cartes interactives basées sur \skill{React.js} et Leaflet.js pour la visualisation.
  \item Conception d'une API REST \skill{Flask} avec \skill{MySQL}, améliorant les performances du tableau de bord grâce à l'optimisation des requêtes et à la mise en cache.
\end{itemize}

% ── Formation ───────────────────────────────────────────────────────────────
\section{Formation}

\entry{Diplôme de niveau Master~: \href{https://www.francecompetences.fr/recherche/rncp/39774}{Architecte Expert en Technologies de l'Information} (RNCP Niveau 7, Bac+5)}{Prévu Févr. 2027}\\
\subentry{\href{https://www.42network.org/}{42}, Paris, France \enspace\textbullet\enspace Spécialisation~: Architecture des Données}{}

% ── Projets ─────────────────────────────────────────────────────────────────
\section{Projets}

\entry{\href{https://github.com/mylastresort/taskmaster}{Système de Contrôle de Tâches}}{Sept. 2026}\\
\subentry{Go, Orchestration de Tâches Asynchrones}{}
\begin{itemize}
  \item Développement d'un système d'orchestration de tâches avec une file interne traitant plus de 1\,000 tâches asynchrones/min et 15+ options de configuration.
  \item Implémentation d'un pool de workers avec arrêt gracieux, files dead-letter et reprises automatiques afin de renforcer la tolérance aux pannes du système.
  \item Construction du cœur client/serveur \skill{Go} avec configuration \skill{TOML} de type superviseur, livré via \skill{Docker Compose} avec tests CI.
\end{itemize}

\vspace{0.40em}

\entry{\href{https://github.com/mylastresort/multilayer-perceptron}{Perceptron Multicouche (MLP)}}{Août 2026}\\
\subentry{Rust, Réseaux de Neurones, Backpropagation from Scratch}{}
\begin{itemize}
  \item Développement d'un MLP configurable avec \skill{rétropropagation} de zéro en \skill{Rust}, piloté par des configurations \skill{YAML}.
  \item Implémentation des optimiseurs \skill{SGD} et \skill{Adam}, early stopping et métriques par époque avec courbes d'apprentissage.
\end{itemize}

\vspace{0.40em}

\entry{\href{https://github.com/mylastresort/hypertube}{Plateforme de Streaming Vidéo}}{Déc. 2025}\\
\subentry{Go, Next.js, PostgreSQL, Node.js, OAuth2}{}
\begin{itemize}
  \item Développement d'un backend \skill{Go} scalable avec streaming HLS et gestion de cache optimisée pour une diffusion vidéo efficace.
  \item Conception d'un schéma \skill{PostgreSQL} avec optimisation des requêtes et pooling de connexions.
  \item Conduite de la livraison \skill{Agile} avec 3 développeurs, via sprint planning et stand-ups quotidiens.
\end{itemize}

% ── Compétences techniques ──────────────────────────────────────────────────
\section{Compétences techniques}

{\small
\renewcommand{\arraystretch}{1.02}
\begin{tabularx}{\linewidth}{@{} >{\bfseries}l X @{}}
  Langages            & TypeScript, JavaScript, Python, Go, Java, C, C++, OCaml, SQL \\
  Frontend            & React, Next.js, Angular, HTML/CSS, Tailwind CSS, Leaflet.js \\
  Backend \& Données  & Node.js, NestJS, FastAPI, Spring Boot, API REST, PostgreSQL, MySQL, MongoDB, Redis, Kafka, FFmpeg \\
  Cloud \& DevOps     & Docker, Docker Compose, Kubernetes, Linux, Vagrant, ArgoCD, GitLab CI/CD, GitHub Actions,\newline Git, AWS, IaC, Terraform, Prometheus, Grafana, Administration Système \\
  Architecture        & Systèmes Distribués, Architecture Événementielle, Microservices, Design Patterns, Tolérance aux Pannes, Scalabilité \\
  IA \& Agents        & Google ADK, MCP, Systèmes Multi-Agents, Ollama (LLM locaux), Tool Calling, Prompt Engineering, Évaluation de LLM \\
  IoT \& Temps Réel   & MQTT, WebSockets, Jumeaux Numériques, Traitement de Données Temps Réel \\
  Sécurité            & OAuth2, JWT, OWASP Top 10 \\
\end{tabularx}
}

\end{document}
